\documentclass[a4paper,12pt]{report}
\usepackage{geometry}
\geometry{a4paper, top=-1cm, bottom=1in, left=0.5in, right=0.5in}
\usepackage{graphicx}
\usepackage[hidelinks]{hyperref}
\usepackage{amsmath}
\usepackage{listings}
\usepackage{bookmark}
\usepackage{xcolor}
\usepackage[utf8]{inputenc}  % Pour supporter les caractères UTF-8
\usepackage[T1]{fontenc}      % Pour une meilleure gestion des caractères accentués
\usepackage{tikz}
\usetikzlibrary{calc}
\usepackage{atbegshi}  % Pour ajouter des éléments sur toutes les pages
\usepackage{fancybox}
\usepackage{mathpazo}
\usepackage{titlesec}
\usepackage{tcolorbox}
\usepackage{caption}
\usepackage{hyperref}
\usepackage{tocloft}
\usepackage{enumitem} % Pour personnaliser les puces
\usepackage{amssymb} % Pour accéder aux symboles de type carré
\usepackage{float}     % Pour utiliser l'option H dans les figures
\usepackage[table]{xcolor} % Nécessaire pour les couleurs des tableaux
\usepackage{tikz}
\usepackage{pifont}   % Pour des symboles supplémentaires
\usepackage{background}



% Configuration par défaut : pas de fond
\backgroundsetup{
  scale=0.5,
  color=black,
  opacity=0.2,
  angle=0,
  vshift=-2cm, % Décalage vertical
  contents={}
}

% Commande pour activer le fond uniquement pour une section précise
\newcommand{\EnableIntroBackground}{
  \backgroundsetup{
    contents={\includegraphics[width=\paperwidth,height=\paperheight]{joli devops.png}}
  }
}

% Commande pour désactiver le fond (par défaut pour toutes les autres pages)
\newcommand{\DisableBackground}{
  \backgroundsetup{
    contents={}
  }
}

\lstdefinelanguage{HCL}{
  morekeywords={resource, provider, variable, output, module, data, terraform},
  sensitive=true,
  morecomment=[l][\color{green}]{\#},
  morestring=[b]"'
}


\titleformat{\chapter}[block]
  {\normalfont\bfseries\Huge} % Police, gras, taille (ajustez 'Huge' selon vos besoins)
  {\begin{tcolorbox}[colback=blue!20, colframe=blue, sharp corners, boxrule=0.8mm, width=1.5cm, halign=center, nobeforeafter] \thechapter \end{tcolorbox}}{1em}{} % Fond coloré autour du numéro

% Personnalisation de l'espacement entre le numéro et le titre
\titleformat{\chapter}[hang]
  {\normalfont\bfseries\Huge}
  {\begin{tcolorbox}[colback=blue!20, colframe=blue, sharp corners, boxrule=0.8mm, width=1.5cm, halign=center, nobeforeafter] \thechapter \end{tcolorbox}}{1em}{}
  [\titlerule] % Ajoute une règle sous le titre du chapitre


% Définir les couleurs et les paramètres pour shadowbox
\setlength{\shadowsize}{3pt} % Taille de l'ombre
\setlength{\fboxsep}{10pt}   % Espace entre le texte et la boîte
\setlength{\fboxrule}{1pt}   % Épaisseur de la bordure
%\colorlet{shadowcolor}{black!40} % Couleur de l'ombre
\colorlet{boxcolor}{blue!70}     % Couleur de la boîte



% Définir une bordure sur toutes les pages
\AtBeginShipout{\AtBeginShipoutAddToBox{%
  \begin{tikzpicture}[remember picture,overlay,inner sep=0,outer sep=0]
    % Bordure extérieure
    \draw[blue!70!black, line width=4pt]
      ([xshift=-0.8cm, yshift=-0.8cm]current page.north east) -- 
      ([xshift=0.8cm, yshift=-0.8cm]current page.north west) -- 
      ([xshift=0.8cm, yshift=0.8cm]current page.south west) -- 
      ([xshift=-0.8cm, yshift=0.8cm]current page.south east) -- cycle;

    % Bordure intérieure (1er niveau)
    \draw[black!50, thin]
      ([xshift=-1.1cm, yshift=-0.5cm]current page.north east) -- 
      ([xshift=1.1cm, yshift=-0.5cm]current page.north west) -- 
      ([xshift=1.1cm, yshift=0.5cm]current page.south west) -- 
      ([xshift=-1.1cm, yshift=0.5cm]current page.south east) -- cycle;

    % Bordure intérieure (2e niveau)
    \draw[black!30, black]
      ([xshift=-0.5cm, yshift=-1.1cm]current page.north east) -- 
      ([xshift=0.5cm, yshift=-1.1cm]current page.north west) -- 
      ([xshift=0.5cm, yshift=1.1cm]current page.south west) -- 
      ([xshift=-0.5cm, yshift=1.1cm]current page.south east) -- cycle;
  \end{tikzpicture}%
}}

% Définir le titre global
\title{
  \begin{minipage}{\textwidth}
    \centering
    \begin{minipage}{0.49\textwidth}
      \centering
      \includegraphics[width=2cm]{UCAD-removebg-preview.png} \par
      {\color{blue}\footnotesize
      \centering
      Université Cheikh Anta Diop de Dakar \vspace{-1em} \newline
      Faculté des Sciences et Techniques \vspace{-1em} \newline
      Département Mathématiques et Informatique}
          \end{minipage}
    \hfill
    \begin{minipage}{0.47\textwidth}
      \centering
      \vspace{-0.3em} % Diminue l'espacement vertical avant ce bloc
      \includegraphics[width=2cm]{image.png} \par
      {\color{blue}\footnotesize 
        Laboratoire d'Algèbre de Cryptographie de\vspace{-1em} \newline 
        Géométrie Algébrique et Applications \vspace{-1em} \newline 
        LACGAA}
    \end{minipage}
  
    \vspace{4em} % Ajout d'un espace après les images alignées
    \par
    {\normalsize  Licence 2, Transmission de Données et Sécurité de l'Information} \par
    \vspace{1em}
    
    \shadowbox{\color{boxcolor}%
      \begin{minipage}{0.9\textwidth}
        \centering
        \textbf{Thème :} \par
        \vspace{0.5em}
        \textbf{\color{black}Surveillance et Alerting des Applications en Production} 
      \end{minipage}
    }
  \end{minipage}
}
\author{
  \\[1em] Présenté et soutenu par \hspace{10em} Sous la direction de
  \\ \textcolor{blue}{M. Mouhamad Al Amine Ba} \hspace{8em} \textcolor{blue}{Mme Diallo Aminata } \\[3em] 
  \fcolorbox{blue}{gray!10}{\begin{minipage}{0.9\textwidth}  % Largeur du cadre ajustée
    \begin{center} % Centrer uniquement le titre "Jury"
      \textcolor{blue}{ \textbf{\huge Jury}} \\[1em] % Titre centré
      \end{center}      \textcolor{blue}{President} :  Dr Demba sow \hspace{4.7cm}  \textcolor{blue}{ UCAD}  \\ 
      \begin{tabbing}
         \hspace{2.3cm} \= \hspace{3cm} \= \kill % Alignement des colonnes
         \textcolor{blue}{Membres} : \> Dr Abdoulaye Thioune \hspace{3cm} \textcolor{blue}{ UCAD} \=  \\[0.5em]  
         \> Cheikh Mbengue \hspace{4.36cm} \textcolor{blue}{ UCAD} \= \\[0.5em] 
         \> Mme Aminata Diallo  \hspace{3.5cm} \textcolor{blue}{ UCAD}
      \end{tabbing}
    \end{minipage}} \\[7em]   
    \textbf{Année académique 2023-2024}
    }


\begin{document}

\date{} % Supprime la date

\maketitle

\newgeometry{top=2cm, bottom=1in, left=1in, right=1in}

\renewcommand{\contentsname}{Table des matières}
\renewcommand{\listfigurename}{Table des figures} % Renomme "List of Figures" en "Table des figures"
\renewcommand{\listtablename}{Table des tableaux} % Renomme "List of Figures" en "Table des figures"


% Personnalisation des puces en étoiles et élégance typographique

\chapter*{Dédicaces}  % Section sans numéro
\addcontentsline{toc}{chapter}{Dédicace}  % Ajouter à la table des matières
{\centering  % Centrage uniquement pour l'image

\includegraphics[width=6cm]{Dedicace.png}

\begin{flushleft}
\textbf{Bismillahi Rahmani Rahim}, au nom de Dieu le Tout Miséricordieux, le Très Miséricordieux.\\[2em]

Je dédie ce mémoire :  
\begin{itemize}[label={\huge\(\star\)}, labelsep=0.5em, left=0pt]
    \item À ma très chère mère \textcolor{blue}{\textbf{Die Mbaye}} et à mon père \textcolor{blue}{\textbf{Al AMine}}.\\
    Puisse Allah vous protéger, vous accorder santé, bonheur et longue vie, et fasse en sorte que je ne vous déçoive jamais.  
    \item À mes chers frères et sœurs : \textcolor{blue}{\textbf{Freere}}, \textcolor{blue}{\textbf{Autres}},.\\ 
    Puisse Allah vous garder, éclairer votre route et vous aider à réaliser vos rêves les plus chers.  
    \item À mes \textcolor{blue}{oncles et tantes}. Que le tout puissant vous accorde une longue vie.  
    \item À mes \textcolor{blue}{amis, frères et camarades de promotion}. Qu'Allah swt nous facilite la suite de notre parcours.  
\end{itemize}
\end{flushleft}

\chapter*{Remerciements}  % Section sans numéro
\addcontentsline{toc}{chapter}{Remerciements}  % Ajouter à la table des matières

\begin{flushleft}
    Louange à Allah le Tout-Puissant, qui m'a guidé et donné la force et la volonté de réaliser ce travail.\textbf{AL HAMDOULILLAH RABBI AL'ALAMINE.} \\[3em]
    Tous ces mots ne sauraient exprimer la gratitude, l'amour, le respect et la reconnaissance. 
    
    \raggedleft
    Je tiens à exprimer ma profonde gratitude à :  
    \begin{itemize}[label={\huge\textbf{$\heartsuit$}}, labelsep=0.5em, left=0pt]  % Augmentation de la taille des puces
        \raggedleft
    
        \item À nos parents  \\ 
        \textcolor{blue}{\textbf{text}} \\ 
        text
        
        \item À notre encadreur \\ 
        \textcolor{blue}{\textbf{Madame Aminata NGOM}} \\ 
        Pour son accompagnement, sa patience et ses encouragements,\\ constants tout au long de ce travail.
       
        \item À nos proches  \\ 
        \textcolor{blue}{\textbf{Mr Dame Kebe}} \\ 
        Pour leur soutiens, son apport tout au long de notre cursus.
        
        \item À mes camarades de classe \\ 
        \textcolor{blue}{\textbf{Moussa Thiam, Rokhyatou Sarr Fall, Dié Mbaye}} \\ 
        Pour leurs conseils, planifications d'études et échanges fructueux, je vous suis reconnaissant.
    \end{itemize}
\end{flushleft}

\vspace{1cm}

\includegraphics[width=15cm]{merci.png}  % L'image n'est pas centrée ici
}


\tableofcontents  % Table des matières
\newpage
\listoffigures % Liste des figures
\newpage
\listoftables % Liste des tables

 %\chapter*{Liste des acronymes}  % Section sans numéro
\addcontentsline{toc}{chapter}{Liste des acronymes}  % Ajouter à la table des matières

\begin{flushleft}
    \begin{itemize}
        \item \textbf{SSH} : Secure Shell  
        \item \textbf{YUM} : Yellowdog Updater
        \item \textbf{CD} : Déploiement Continu  
        \item \textbf{AWS} : Amazon Web Services  
        \item \textbf{CI} : Intégration Continue  
        \item \textbf{TSDB} : Time Series Database
        \item \textbf{CLI} : Command Line Interface 
        \item \textbf{RPM} : Red Hat Package Manager 
        \item \textbf{CPU} : Central Processing Unit  
        \item \textbf{URL} : Uniform Resource Locator
        \item \textbf{HTTP} : HyperText Transfer Protocol, 
        \item \textbf{TCP} : Transmission Control Protocol
        \item \textbf{Amazon EC2} : Amazon Elastic Compute Cloud  
        \item \textbf{HTTPS} : HyperText Transfer Protocol Secure
        \item \textbf{API} : Interfaces de programmation d'applications  

    \end{itemize}
    \end{flushleft}
    
    % Configuration de l'image de fond

    \chapter*{Introduction}  % Section sans numéro
    \EnableIntroBackground % Activer le fond
    \addcontentsline{toc}{chapter}{Introduction}  % Ajouter à la table des matières
    « Plutôt qu'un marché, DevOps est avant tout une philosophie, un changement culturel qui associe les mondes du développement et des opérations », explique Laurie Wurster, directrice de recherche chez Gartner. Depuis sa naissance en 2009, le mouvement DevOps a révolutionné la façon dont les entreprises conçoivent, développent et déploient leurs applications.\\

    Des entreprises emblématiques de la Silicon Valley, telles que Google, Microsoft, Amazon, Netflix et LinkedIn, ont rapidement adopté DevOps pour améliorer leurs performances, stimuler leur capacité d’innovation et renforcer leur compétitivité. \\
    
    Dans cet écosystème, des outils clés jouent un rôle central pour automatiser, surveiller et gérer les performances des applications dans des environnements complexes et dynamiques. Ces technologies facilitent non seulement le déploiement des applications, mais elles garantissent aussi leur disponibilité, leur scalabilité et la qualité du service.\\
    
    L’intégration de ces outils dans une infrastructure moderne assure une gestion optimale des applications tout en offrant des capacités avancées de surveillance. Cela permet une réactivité accrue face aux problèmes de performance, tout en soutenant la stabilité des systèmes et en favorisant l'innovation et la croissance des entreprises. \\
    
    Prometheus, par exemple, interroge régulièrement les points de terminaison d’une application pour collecter des métriques essentielles à son bon fonctionnement. Cependant, la collecte des métriques n’est qu’une première étape. Leur interprétation et leur visualisation sont tout aussi importantes. 
    
    Grafana joue ici un rôle crucial en tant que tableau de bord visuel, permettant de surveiller, analyser et présenter ces métriques de manière intuitive. Sa flexibilité et ses nombreuses options de personnalisation en font un outil incontournable pour les équipes DevOps.
    
    Enfin, la gestion des alertes via Alertmanager simplifie la supervision en connectant les alertes aux canaux de communication tels que les emails ou Slack. Cette automatisation renforce la réactivité des administrateurs, qui peuvent agir rapidement pour résoudre les problèmes et minimiser les interruptions de service.
    


\DisableBackground


 %\include{chap1}

% Exemple : Police personnalisée à 30pt
% Espacement pour centrer le texte verticalement et horizontalement
 %\include{chap2}
 %\include{chap3}
%\include{End}

\end{document}